\chapter{Background and Literature Review}
\label{ch:lit}

\section{Life Cycle Assessment}

\subsection{The four phases}

ISO~14040/44 define LCA as an iterative four-phase procedure \citep{iso14040,iso14044}
(Figure~\ref{fig:lca-phases}):

\begin{enumerate}
  \item \textbf{Goal and scope definition} -- purpose, intended audience, the \emph{functional
  unit} (here, 1~kg of food as sold/consumed) and the \emph{system boundary} (here,
  cradle-to-consumption: agriculture, processing, packaging, transport, retail, consumer use).
  \item \textbf{Life Cycle Inventory (LCI)} -- compilation of all input and output flows
  (materials, energy, emissions, land, water) per functional unit.
  \item \textbf{Life Cycle Impact Assessment (LCIA)} -- conversion of flows into impact indicators
  by multiplying each flow by a \emph{characterisation factor}; e.g.\ for climate change,
  $\text{GWP} = \sum_i m_i \cdot \text{CF}_i$ in kg~\co{}, where $m_i$ is the mass of greenhouse
  gas $i$ emitted and $\text{CF}_i$ its 100-year global warming potential.
  \item \textbf{Interpretation} -- hotspot identification, sensitivity and uncertainty analysis,
  conclusions and recommendations.
\end{enumerate}

\begin{figure}[ht]
\centering
\resizebox{\textwidth}{!}{%
\begin{tikzpicture}[node distance=6mm and 10mm, every node/.style={font=\small},
  box/.style={draw=c1, fill=c1!8, rounded corners=3pt, minimum width=3.4cm, minimum height=1cm, align=center},
  ml/.style={draw=c2, fill=c2!10, rounded corners=3pt, align=center, font=\footnotesize, text width=4.2cm}]
\node[box] (g) {Goal \& scope\\definition};
\node[box, below=of g] (i) {Inventory analysis\\(LCI)};
\node[box, below=of i] (a) {Impact assessment\\(LCIA)};
\node[box, right=22mm of i, minimum height=4.2cm, minimum width=2.6cm] (int) {Interpretation};
\draw[-{Latex}, thick, c1] (g) -- (i); \draw[-{Latex}, thick, c1] (i) -- (a);
\foreach \n in {g,i,a} \draw[{Latex}-{Latex}, c1] (\n.east) -- (\n.east -| int.west);
\node[ml, left=10mm of i] (m1) {\textbf{ML here:} gap filling, inventory mapping (\textbf{Module B})};
\node[ml, left=10mm of a] (m2) {\textbf{ML here:} surrogate impact prediction (\textbf{Module A}), CF prediction};
\node[ml, right=6mm of int] (m3) {\textbf{ML here:} explanation (SHAP), uncertainty (conformal)};
\draw[-{Latex}, dashed, c2] (m1) -- (i); \draw[-{Latex}, dashed, c2] (m2) -- (a); \draw[-{Latex}, dashed, c2] (m3) -- (int);
\end{tikzpicture}}
\caption{The four ISO 14040/44 phases of LCA and where AI/ML methods integrate (dashed). Bold labels mark the parts implemented in this thesis.}
\label{fig:lca-phases}
\end{figure}

\subsection{Impact categories and the EF 3.0 method}

Results in this thesis use the European Commission's Environmental Footprint (EF~3.0) LCIA method
\citep{fazio2018}, which reports 16 midpoint indicators -- climate change, ozone depletion,
ionising radiation, photochemical ozone formation, particulate matter, human toxicity (cancer and
non-cancer), acidification, eutrophication (freshwater, marine, terrestrial), freshwater
ecotoxicity, land use, and depletion of water, fossil and mineral resources -- and a weighted
\emph{single score} in milli-points (mPt). Many of these indicators are strongly correlated across
products \citep{steinmann2016}, which suggests that a model which learns one indicator well may
transfer to others; this is tested in Section~\ref{sec:e3}.

\subsection{Types of LCA data}

\begin{itemize}
  \item \textbf{Process-based LCI databases} (e.g.\ ecoinvent \citep{wernet2016}, AGRIBALYSE
  \citep{agribalyse2020}) model individual unit processes in physical units; accurate but costly
  and (in ecoinvent's case) licensed.
  \item \textbf{Environmentally extended input--output (EEIO) models} (e.g.\ USEEIO
  \citep{yang2017useeio}) give impacts per unit of monetary spend by economic sector; complete
  coverage but coarse resolution.
  \item \textbf{Product-level result databases} such as the AGRIBALYSE synthesis tables give
  ready-computed impacts for finished products, which is precisely the supervised-learning target
  needed here.
\end{itemize}

\subsection{Why LCA is slow: the data bottleneck}

Reap et al.\ \citep{reap2008} identified data availability and quality as a central unresolved
problem in LCA, and Hellweg and Milà~i~Canals \citep{hellweg2014} highlight the need for faster,
streamlined approaches that keep scientific robustness. Two specific symptoms motivate this work:
(i) \emph{missing data} for new or niche products, and (ii) the \emph{mapping effort} needed to
connect product descriptions to database processes.

\section{AI/ML in LCA: a structured review}

Ghoroghi et al.\ \citep{ghoroghi2022} systematically reviewed ML applications in LCA and found
them concentrated in optimisation, prediction of impacts, and support for inventory compilation,
while noting gaps in transparency, uncertainty treatment and generalisability. Algren et al.\
\citep{algren2021} similarly survey ML methods for handling incomplete or uncertain LCA data.
The literature can be organised by the LCA phase that ML supports (Table~\ref{tab:lit}).

\subsection{Inventory gap filling}

Hou et al.\ \citep{hou2018} estimated missing unit-process flows in ecoinvent with a
similarity-based approach, accurate when fewer than about 5\% of flows are missing. Zhao et al.\
\citep{zhao2021} extended this with decision trees, classifying zero/non-zero flows and estimating
non-zero values with $R^2>0.7$ when under 20\% of a process is missing. These works operate on the
inventory matrix itself and require licensed process data.

\subsection{Surrogate models for rapid impact screening}

Song, Keller and Suh \citep{song2017} trained neural networks on molecular descriptors to predict
the cradle-to-gate impacts of chemicals, reaching $R^2$ of 0.87 for the Eco-indicator 99 score but
only 0.48 for global warming, and stressed that \emph{applicability domains} matter when
interpreting such screening predictions. This is the closest conceptual precedent for Module~A:
learn a direct map from cheap descriptors to impacts. In the food domain, Clark et al.\
\citep{clark2022} estimated four environmental indicators for 57{,}000 UK and Irish products by
inferring ingredient composition from ingredient lists and pairing it with LCA data -- a
rule-based rather than learned approach, but one that shows the value of ingredient information.

\subsection{NLP for mapping products to LCA data}

Balaji et al.\ \citep{balaji2023caml} (CaML) used pre-trained sentence-transformer embeddings
\citep{reimers2019} to map product descriptions to EEIO industry sectors in a zero-shot manner,
ranking the top-5 candidates for human confirmation and reporting a mean absolute percentage error
of about 22\% without domain labels. This motivates Module~B, which maps to \emph{product-level
process datasets} instead of sectors, compares several encoders, and tests a cross-lingual setting.

\subsection{Explainability and uncertainty}

Tree ensembles such as Random Forests \citep{breiman2001} and gradient boosting (XGBoost
\citep{chen2016xgboost}) are strong on tabular data, and SHAP values \citep{lundberg2017,lundberg2020}
provide additive, locally exact explanations for them. For uncertainty, conformal prediction
\citep{vovk2005,lei2018,angelopoulos2023} yields prediction intervals with finite-sample coverage
guarantees that do not depend on the model being correct -- an attractive property for LCA, where
uncertainty reporting is required by ISO~14044 but is rarely provided by ML surrogates.

\begin{table}[ht]
\centering\small
\caption{Representative AI/ML-in-LCA studies and how this thesis relates to them.}
\label{tab:lit}
\begin{tabularx}{\textwidth}{@{}p{2.6cm}p{2.4cm}Xp{2.3cm}@{}}
\toprule
Study & LCA phase & Approach & Gap addressed here \\
\midrule
Hou et al.\ 2018 \citep{hou2018} & Inventory & Similarity-based estimation of missing flows (ecoinvent) & Open data; product level \\
Zhao et al.\ 2021 \citep{zhao2021} & Inventory & Decision trees for missing flows & Uncertainty intervals \\
Song et al.\ 2017 \citep{song2017} & Impact & ANN from molecular descriptors to impacts & Explainability; unseen-category test \\
Clark et al.\ 2022 \citep{clark2022} & Inventory/Impact & Ingredient-list inference for 57k foods & Learned model; uncertainty \\
Balaji et al.\ 2023 \citep{balaji2023caml} & Inventory mapping & Zero-shot sentence similarity to EEIO sectors & Process-level, multilingual, encoder benchmark \\
Ghoroghi et al.\ 2022 \citep{ghoroghi2022} & Review & Survey of ML in LCA & -- \\
\midrule
\textbf{This thesis} & Mapping + Impact + Interpretation & XGBoost surrogate + E5 mapping + SHAP + conformal, integrated tool & -- \\
\bottomrule
\end{tabularx}
\end{table}

\section{Research gap}

From the review, four gaps emerge that this thesis addresses together:
\begin{enumerate}
  \item \textbf{Evaluation realism.} Most surrogate studies report random-split accuracy, which
  rewards memorising near-duplicates. Performance on genuinely new categories is rarely reported.
  \item \textbf{Uncertainty.} ML predictions in LCA are usually point estimates without calibrated
  intervals.
  \item \textbf{Transparency.} Black-box predictions are hard for LCA practitioners to trust or audit.
  \item \textbf{Integration.} Estimation and mapping are studied separately, though in practice
  they form one workflow (describe $\rightarrow$ map $\rightarrow$ estimate $\rightarrow$
  interpret).
\end{enumerate}
